\documentclass[11pt,a4paper]{article}

\usepackage[margin=1in]{geometry}
\usepackage{amsmath,amssymb}
\usepackage{booktabs}
\usepackage{graphicx}
\usepackage{hyperref}
\usepackage{natbib}
\usepackage{microtype}
\usepackage{xcolor}
\usepackage{multirow}
\usepackage{array}
\usepackage{caption}

\hypersetup{colorlinks=true, linkcolor=blue, citecolor=blue, urlcolor=blue}

\title{Probing and Steering Uncertainty in Biomedical Language Models}
\author{Debmalya Pal}
\date{}

\begin{document}
\maketitle

\begin{abstract}
Biomedical language models are increasingly deployed in clinical decision
support systems, yet tend to generate assertions in a confident register
even when clinical evidence is incomplete. We investigate whether
\emph{uncertainty}---the hedged epistemic stance that clinical
professionals express through phrases like ``findings are consistent
with'' or ``cannot exclude''---is linearly represented in the hidden
states of biomedical language models and can be controlled through
activation steering at inference time without retraining. We study seven
models spanning two architectural families: three decoder (causal LM)
models---BioGPT \citep{luo2022biogpt}, BioGPT-Large, and BioMedLM---and
four encoder (masked LM) models---BioBERT \citep{lee2020biobert},
ClinicalBERT \citep{alsentzer2019clinicalbert}, BlueBERT
\citep{peng2019bluebert}, and SciBERT \citep{beltagy2019scibert}. Using
BioScope speculation annotations \citep{vincze2008bioscope} as ground
truth, we show that uncertainty is linearly decodable across all seven
models (5-fold CV accuracy 83--87\%), with probing curves that are
uniformly strong across middle-to-late layers. For decoder models, we
demonstrate that activation steering can increase hedge scores at low
intervention magnitudes, but the productive range is narrow and
generation quality degrades above a model-specific threshold. For encoder
models, we quantify the representation shift induced by steering along
the probe direction using five complementary metrics, finding consistent,
measurable shifts even at small magnitudes. We introduce a principled
layer-selection procedure combining a 2/3-depth prior with local
empirical refinement, and show it is empirically validated by probe
accuracy data. We close with a methodological argument that multi-metric
evaluation is necessary for steering interventions in safety-sensitive
domains.
\end{abstract}

%% -----------------------------------------------------------------------
\section{Introduction}
%% -----------------------------------------------------------------------

Biomedical language models are increasingly proposed as components of
clinical decision support systems, from radiology report generation to
literature-assisted diagnosis. BioGPT \citep{luo2022biogpt}, trained on
15M PubMed abstracts, has become a commonly used open-source model for
these tasks, achieving strong performance on biomedical question
answering and relation extraction benchmarks. Yet a gap persists between
benchmark performance and clinical deployability: language models trained
on biomedical text tend to produce assertions in a confident declarative
register, even when the clinical evidence is incomplete, ambiguous, or
insufficient to warrant such certainty \citep{singhal2023medpalm,
umapathi2023medhalt}.

This overconfidence is not merely a stylistic issue. In clinical
practice, expressing appropriate uncertainty---through hedged language
such as ``findings are consistent with,'' ``cannot exclude,'' or ``may
represent''---is a learned professional skill with real safety
consequences \citep{panicek2016radiology}. A downstream user who reads a
confident model output is more likely to act on it without seeking
confirmation, whereas a user who reads a hedged output is prompted to
consider alternatives. If biomedical language models are to participate
in clinical workflows, controlling the expression of uncertainty in
their outputs is therefore not optional but central.

The dominant approaches to controlling model behavior fall into two
families. The first operates at the input level through prompting ---
instructing the model to be cautious or to hedge. This approach is cheap
and requires no access to model internals, but its effects are known to
be brittle \citep{zhao2021calibrate, lu2022order}. The second family
operates at the parameter level through fine-tuning, instruction tuning,
or RLHF. These methods can produce more reliable behavioral changes but
require substantial labeled data, compute, and engineering effort.

A third approach has emerged from mechanistic interpretability: directly
intervening on a model's hidden states at inference time, along
directions in activation space that correspond to target concepts. This
family of methods---\emph{activation steering}
\citep{turner2023activation}, \emph{representation engineering}
\citep{zou2023representation}, or \emph{contrastive activation addition}
\citep{panickssery2024steering}---rests on the empirical claim that many
high-level concepts are linearly represented in transformer hidden
states, and that adding a scaled direction vector at inference can
predictably shift model behavior.

The premise of our work is to bring this methodology to bear on clinical
uncertainty across a representative set of biomedical language models.
The appeal is clear: if uncertainty were a steerable direction, we could
modulate a model's hedging behavior at inference time without retraining.
Whether this prospect is realistic is an empirical question that has not
previously been investigated across the biomedical model landscape.

Two conditions must hold. First, uncertainty must be represented in the
hidden states as a structure that linear intervention can act on. Second,
when we intervene along such a direction, the model's output must shift
in the target direction without collapsing other properties we care
about. The first is a question about representational structure; the
second is a question about the causal consequences of intervention.

Our contribution is to investigate both questions systematically across
seven biomedical language models, using BioScope-annotated speculation
labels \citep{vincze2008bioscope} as ground truth and a combination of
linear probing and activation steering as our methodological lens. We
make five specific contributions:

\begin{itemize}
\item We demonstrate that uncertainty is linearly decodable across all
  seven tested models (decoder and encoder alike), with 5-fold CV
  accuracies of 83--87\%, establishing that linear encoding of epistemic
  stance is a general property of biomedical language models rather than
  a model-specific artifact.

\item We show that the choice of representation extraction method has a
  dramatic effect on observable linear structure in BioGPT (Cohen's
  $d = 1.33$ vs.\ $0.09$ for mean vs.\ last-token pooling), with
  methodological implications for probing biomedical LMs broadly.

\item We introduce a principled two-stage layer selection procedure
  combining a 2/3-depth prior with local empirical refinement (best
  accuracy within $\pm 2$ layers), and validate it empirically: for five
  of seven models, the window agrees with the strict 2/3 anchor; for
  ClinicalBERT, the window correctly identifies a substantially better
  layer ($+3.25$pp) that the naive anchor misses.

\item We characterize decoder steering using four output metrics,
  demonstrating a distinct productive range at low magnitudes and a
  collapse regime at larger magnitudes, and show that dynamic hidden-norm
  calibration substantially extends the productive range compared to
  static norm estimates.

\item We introduce an encoder representation-steering evaluation suite
  (five metrics: projection shift, probability shift, decision-boundary
  flip rate, centroid cosine shift, and on-axis fraction) and show
  consistent, measurable representation shifts across all four encoder
  models, even at small steering magnitudes.
\end{itemize}

We intend for this paper to be useful in two ways: as positive evidence
that biomedical language models encode semantic properties in
surprisingly linear ways, and as a cautionary case study showing that
this structural finding does not directly translate into clean behavioral
control. Both findings should inform how the field approaches
inference-time intervention in clinical language models.

%% -----------------------------------------------------------------------
\section{Related Work}
%% -----------------------------------------------------------------------

Our work sits at the intersection of three active research areas:
uncertainty in biomedical and clinical language processing, probing of
pretrained language model representations, and activation-level control
of generative language models.

\paragraph{Uncertainty and hedging in biomedical text.}
The linguistic study of hedging and speculation in biomedical writing has
a long history predating modern language models. The BioScope corpus
\citep{vincze2008bioscope} remains the standard annotated resource,
marking speculation cues and their scope in abstracts, full papers, and
clinical notes. Associated work established classification benchmarks for
sentence-level and token-level speculation detection
\citep{agarwal2010detecting, adel2015comparing}, demonstrating that
hedging in biomedical text can be detected with high accuracy using
lexical and contextual features.

Separately, the clinical NLP literature has examined the role of hedging
in specific medical genres, particularly radiology reports, where
cautious phrasing is an explicit professional norm
\citep{panicek2016radiology}. Recent interest in the calibration of
large language models in clinical contexts \citep{umapathi2023medhalt,
singhal2023medpalm} has highlighted that biomedical LMs can produce
plausible-sounding but unsupported clinical claims, though most such work
focuses on factual accuracy rather than the expression of uncertainty per
se. What is largely absent from this literature is the investigation of
uncertainty as an \emph{internal representational property} of
biomedical language models. Our work complements these perspectives by
asking where and how uncertainty is represented inside the model's hidden
states.

\paragraph{Probing pretrained language model representations.}
The probing paradigm---training simple classifiers on frozen hidden
states to investigate what a language model has encoded---has been a
central tool in interpretability research \citep{hewitt2019structural,
tenney2019bert, belinkov2022probing}. Core findings have shown that
language models encode a layered hierarchy of linguistic information:
surface features in early layers, syntactic structure in middle layers,
and semantic properties in later layers. In the biomedical domain,
probing has been applied more narrowly, primarily to investigate what
domain knowledge language models have acquired during pretraining
\citep{sung2021biomedical, vulic2020probing}. Less work has examined
pragmatic or stance-like properties of biomedical LMs. A recurring
concern in the probing literature is the interpretation problem: a
probe's accuracy tells us the information is \emph{decodable} from the
representation, but not necessarily that the model uses this information
in its generative behavior \citep{ravichander2021probing}.

\paragraph{Activation-level intervention in language models.}
The broader research program of intervening directly on language model
activations to control behavior has developed rapidly. The foundational
claim---the \emph{linear representation hypothesis}---is that high-level
concepts in language models are encoded as directions in activation
space, such that modifying these directions at inference time
predictably modulates corresponding behavior \citep{park2023linear}.
Concrete methods include activation addition \citep{turner2023activation},
representation engineering \citep{zou2023representation}, and contrastive
activation addition \citep{panickssery2024steering}. These methods have
been demonstrated for target behaviors including sentiment, refusal, and
truthfulness, primarily in general-purpose language models.
Methodological concerns are increasingly acknowledged: steering vectors
can be confounded with surface features \citep{elazar2021amnesic}, and
aggressive steering can push activations outside the training
distribution \citep{zou2023representation}. In the biomedical setting,
application of these methods is, to our knowledge, essentially
unexplored.

\paragraph{The gap this paper addresses.}
Bringing these threads together, we identify a specific and tractable
gap. The uncertainty literature in biomedical NLP has treated hedging as
an external linguistic property. The probing literature has established
that language models encode rich layered representations but has devoted
less attention to biomedical models and pragmatic features. The
activation steering literature has demonstrated control over
general-domain behaviors but has not been systematically evaluated in
biomedical contexts with multi-metric evaluation, and not across
architecturally diverse biomedical models. Our work takes the
question---can uncertainty in biomedical language models be understood as
a linearly represented and steerable concept?---and investigates it
across seven models, two architectures, and both generation and
representation modalities.

%% -----------------------------------------------------------------------
\section{Methods}
%% -----------------------------------------------------------------------

We describe our methodological pipeline. Where design choices were shaped
by preliminary observations, we note the empirical motivation.

\subsection{Dataset construction}

We use the BioScope corpus \citep{vincze2008bioscope}, which provides
sentence-level annotations of speculation, negation, and their scopes in
biomedical text. A BioScope sentence is labeled as containing speculation
if any of its tokens are marked with a speculation cue
(e.g.\ \texttt{<cue type="speculation">}). We treat this annotation as a
binary sentence-level label: sentences containing any speculation cue are
\emph{uncertain}; sentences containing no speculation cue are
\emph{certain}.

Parsing all three BioScope sub-corpora (abstracts, full papers, clinical
reports) yields 3,065 uncertain and 13,007 certain sentences. To
mitigate confounds between speculation status and trivial surface
properties, we restrict to sentences between 30 and 200 characters in
length, reducing the pool to 1,825 uncertain and 9,810 certain
sentences. From this filtered pool, we sample 200 sentences of each
class (seed 42) to form our balanced contrast set
$\mathcal{D} = \mathcal{D}_u \cup \mathcal{D}_c$ with
$|\mathcal{D}_u| = |\mathcal{D}_c| = 200$. File enumeration during
parsing is alphabetically sorted to ensure cross-platform reproducibility
of the sampled sentences.

Despite length filtering, a small systematic difference remains: uncertain
sentences average 27.1 tokens versus 24.6 for certain sentences---a
difference we later account for through orthogonalization.

\subsection{Models and representation extraction}

We study seven biomedical language models across two architectural
families, summarized in Table~\ref{tab:models}.

\begin{table}[ht]
\centering
\caption{Models evaluated. All experiments use float32 precision;
  quantization was not applied to any model in scope.}
\label{tab:models}
\begin{tabular}{llrrrl}
\toprule
\textbf{Model} & \textbf{Architecture} & \textbf{Layers} & \textbf{Dim} & \textbf{Params} & \textbf{Pretraining corpus} \\
\midrule
BioGPT       & Decoder (causal LM) & 24 & 1024 & 345M  & PubMed abstracts \\
BioGPT-Large & Decoder (causal LM) & 48 & 1600 & 1.5B  & PubMed abstracts \\
BioMedLM     & Decoder (causal LM) & 32 & 2560 & 2.7B  & PubMed abstracts \\
\midrule
BioBERT      & Encoder (masked LM) & 12 & 768  & 110M  & PubMed + PMC \\
ClinicalBERT & Encoder (masked LM) & 12 & 768  & 110M  & MIMIC-III notes \\
BlueBERT     & Encoder (masked LM) & 12 & 768  & 110M  & PubMed + MIMIC-III \\
SciBERT      & Encoder (masked LM) & 12 & 768  & 110M  & Broad science \\
\bottomrule
\end{tabular}
\end{table}

For any input sentence $s$ tokenized into $T$ tokens, we extract the
hidden state tensor $\mathbf{H}_\ell(s) \in \mathbb{R}^{T \times d}$ at
transformer layer $\ell$, then mean-pool over non-padding tokens:
\begin{equation}
\mathbf{h}_\ell(s) = \frac{1}{T}\sum_{t=1}^{T} \mathbf{H}_\ell(s)[t,:]
\label{eq:meanpool}
\end{equation}
As we establish in Section~\ref{sec:pooling} for BioGPT, mean pooling
yields dramatically better representational structure than last-token
pooling (Cohen's $d = 1.33$ vs.\ $0.09$), so we use mean pooling
uniformly across all models.

\subsection{Linear probing}

To characterize where and how strongly uncertainty is linearly encoded
across each model's depth, we train logistic regression probes on
mean-pooled hidden states. For each layer $\ell$:
\begin{equation}
p(y = \text{uncertain} \mid \mathbf{h}_\ell(s)) =
\sigma\!\left(\mathbf{w}_\ell^\top \mathbf{h}_\ell(s) + b_\ell\right)
\label{eq:probe}
\end{equation}
with $\ell_2$-regularization via $C = 0.1$ (fixed without
cross-validated tuning). We evaluate each probe under stratified 5-fold
cross-validation on the balanced 400-sentence contrast set, reporting
mean and standard deviation of fold accuracies and a 95\% CI as
$\bar{a} \pm 1.96 \cdot s/\sqrt{k}$ with $k=5$.

\subsection{Principled steering layer selection}
\label{sec:layerselect}

A crucial design choice is which layer to steer. We propose a two-stage
procedure that combines a theoretically motivated prior with local
empirical refinement.

\paragraph{Stage 1: 2/3-depth anchor.}
We define an anchor layer as the layer index closest to $2/3$ of the
model's depth:
\begin{equation}
\ell^* = \mathrm{round}\!\left(\frac{2N}{3}\right)
\label{eq:anchor}
\end{equation}
where $N$ is the number of transformer blocks. The rationale is that
representations at this depth balance two competing objectives:
(i)~\emph{representational maturity}---by $2/3$ depth, hidden states
encode abstract, distributed semantic properties rather than surface
lexical patterns; and (ii)~\emph{propagation distance}---enough
downstream blocks remain for the injected signal to shape output
behavior. For BioGPT (24 layers), this gives $\ell^* = 16$, matching
the empirically identified layer in \citet{luo2022biogpt}.

\paragraph{Stage 2: Local empirical refinement.}
Probe accuracy curves are often locally flat: a plateau spanning several
adjacent layers, all with statistically indistinguishable accuracy. In
such settings, the strict 2/3 anchor may not be the best available
choice within its neighbourhood. We therefore search within a window of
$\pm 2$ layers around the anchor, selecting the layer with the highest
mean probe accuracy, with ties broken by minimum standard deviation (the
more reliable probe across CV folds):
\begin{equation}
\hat{\ell} = \mathop{\mathrm{argmax}}_{\ell \,\in\,
  [\ell^* - 2,\, \ell^* + 2] \,\cap\, \{1,\ldots,N\}}
  \;\bigl(\bar{a}_\ell,\; -s_\ell\bigr)
\label{eq:windowselect}
\end{equation}
where $\bar{a}_\ell$ and $s_\ell$ are the mean and standard deviation
of probe CV accuracy at layer $\ell$. Layer 0 (the embedding output) is
excluded because it does not correspond to a hookable transformer-block
module.

\paragraph{Empirical validation.}
Table~\ref{tab:layerselect} shows the anchor and window-selected layers
for all five models with complete probing results. The window agrees with
the anchor for BioGPT and BioBERT, confirming the heuristic. For
ClinicalBERT, the window selects L6 over the L8 anchor with a
$+3.25$pp accuracy advantage and lower CV variance
($\sigma = 2.98\%$ vs.\ $4.45\%$), a difference that is clearly
meaningful given the within-window range of only 13pp. For BlueBERT and
SciBERT, the window makes a 1-layer adjustment with a 0.25pp accuracy
gain---within noise, but strictly non-negative.

\begin{table}[ht]
\centering
\caption{Steering layer selection. ``Anchor'' is $\ell^* = \mathrm{round}(2N/3)$.
  ``Window best'' is $\hat{\ell}$ per Eq.~\eqref{eq:windowselect}.
  Shaded rows indicate cases where window selection differs from anchor.}
\label{tab:layerselect}
\begin{tabular}{lrrrrrr}
\toprule
\textbf{Model} & $N$ & \textbf{Anchor} $\ell^*$ &
  \textbf{Acc@anchor} & \textbf{Window best} $\hat{\ell}$ &
  \textbf{Acc@}$\hat{\ell}$ & \textbf{Std@}$\hat{\ell}$ \\
\midrule
BioGPT       & 24 & 16 & 85.50\% & 16 & 85.50\% & 2.44\% \\
BioBERT      & 12 &  8 & 86.00\% &  8 & 86.00\% & 3.24\% \\
\rowcolor{gray!15}
ClinicalBERT & 12 &  8 & 81.50\% &  6 & 84.75\% & 2.98\% \\
\rowcolor{gray!15}
BlueBERT     & 12 &  8 & 83.25\% & 10 & 83.50\% & 5.18\% \\
\rowcolor{gray!15}
SciBERT      & 12 &  8 & 86.00\% & 10 & 86.25\% & 4.76\% \\
\bottomrule
\end{tabular}
\end{table}

\subsection{Hidden-state norm calibration}

The steering magnitude is parameterized fractionally:
\begin{equation}
\alpha_{\text{eff}} = \alpha_{\text{frac}} \cdot
\|\mathbf{h}_{\hat{\ell}}\|_{\text{est}}
\label{eq:alphaeff}
\end{equation}
where $\|\mathbf{h}_{\hat{\ell}}\|_{\text{est}}$ is the empirical mean
hidden-state norm at the selected layer. We estimate this dynamically
from 8 generic biomedical sentences, computing per-token L2 norms and
averaging. This makes $\alpha_{\mathrm{frac}}$ interpretable as the
fraction of the hidden state's typical magnitude being perturbed,
independently of which layer or model is being steered.

For BioGPT at layer 16, dynamic calibration yields a norm estimate of
approximately 205, compared to the hard-coded value of 616 used in
earlier experiments on the same model. This factor-of-3 discrepancy had
a concrete consequence: every $\alpha_{\mathrm{frac}}$ step was
effectively 3$\times$ more aggressive than intended, explaining why
previously reported results showed perplexity collapse ($>400$) at
$\alpha_{\mathrm{frac}} = 0.25$. With calibrated norms, the same
$\alpha_{\mathrm{frac}} = 0.25$ corresponds to a much smaller absolute
perturbation, and perplexity remains controlled.

\subsection{Steering direction construction}

We construct two steering directions at the selected layer $\hat{\ell}$.

\paragraph{Probe direction.}
The weight vector of a logistic regression probe trained on the full
contrast set:
\begin{equation}
\mathbf{v}_{\text{probe}} = \mathbf{w}^{\text{full}}_{\hat{\ell}}
\end{equation}

\paragraph{Orthogonalized direction.}
To reduce the contribution of surface confounds, we orthogonalize
$\mathbf{v}_{\text{probe}}$ against a length direction
$\mathbf{v}_{\mathrm{len}}$ (from linear regression predicting token
count) and a hedge-count direction $\mathbf{v}_{\mathrm{hedge}}$ (from
linear regression predicting occurrences of hedge vocabulary):
\begin{equation}
\mathbf{v}_{\text{ortho}} = \mathrm{orth}\!\left(
\mathrm{orth}(\mathbf{v}_{\text{probe}},
\mathbf{v}_{\text{len}}),\,
\mathbf{v}_{\text{hedge}}\right), \quad
\mathrm{orth}(\mathbf{u}, \mathbf{v}) = \mathbf{u} -
\frac{\mathbf{u}^\top\mathbf{v}}{\|\mathbf{v}\|^2}\mathbf{v}
\end{equation}

\subsection{Decoder activation steering}

For decoder models (BioGPT, BioGPT-Large, BioMedLM), we modify the
forward pass via a forward hook registered on layer $\hat{\ell}$:
\begin{equation}
\mathbf{h}^{(t)}_{\hat{\ell}} \leftarrow
\mathbf{h}^{(t)}_{\hat{\ell}} + \alpha_{\text{eff}} \cdot \hat{\mathbf{v}},
\quad \forall t
\end{equation}
where $\hat{\mathbf{v}} = \mathbf{v}/\|\mathbf{v}\|$. The hook fires
continuously during both the prefill phase and each autoregressive
decoding step.

We sweep $\alpha_{\mathrm{frac}}$ across nine values
$\{0, 0.025, 0.05, 0.075, 0.10, 0.125, 0.15, 0.20, 0.25\}$ for both
the probe and orthogonalized directions. Generations use temperature
$T = 0.8$, top-$p = 0.9$, repetition penalty 1.2, minimum 20 and
maximum 60 new tokens. We use five clinical prompt templates covering
chest imaging findings and use 20 random seeds per configuration (100
samples per direction-$\alpha_{\mathrm{frac}}$ pair).

\subsection{Encoder representation steering}

For encoder models (BioBERT, ClinicalBERT, BlueBERT, SciBERT), text
generation is not available. We instead measure the representation-level
effect of steering by injecting the steering vector at $\hat{\ell}$ and
comparing the modified and original mean-pooled hidden states.

For each test sentence $s$ (drawn from a held-out 80-sentence balanced
set not used to train the probe), we compute five metrics:

\begin{enumerate}
\item \textbf{Probe projection shift} ($\Delta_{\text{proj}}$): change in
  scalar projection onto the unit probe direction,
  $\Delta_{\text{proj}} = (\mathbf{h}_{\hat{\ell}}^{\text{mod}} -
  \mathbf{h}_{\hat{\ell}}^{\text{orig}}) \cdot \hat{\mathbf{v}}_{\text{probe}}$.

\item \textbf{Probability shift} ($\Delta_{\text{prob}}$): change in the
  probe classifier's predicted probability of the ``uncertain'' class,
  $\Delta_{\text{prob}} = p(y{=}1 \mid \mathbf{h}^{\text{mod}}) -
  p(y{=}1 \mid \mathbf{h}^{\text{orig}})$.

\item \textbf{Decision-boundary flip rate} ($\text{flip}\%$): fraction of
  sentences where the probe's binary prediction changes from certain to
  uncertain.

\item \textbf{Centroid cosine shift} ($\Delta_{\text{cos}}$): change in
  cosine similarity between the sentence representation and the mean
  hidden state of the uncertain training sentences.

\item \textbf{On-axis fraction}: the fraction of the total hidden-state
  shift $\Delta\mathbf{h} = \mathbf{h}^{\text{mod}} - \mathbf{h}^{\text{orig}}$
  that aligns with the probe direction,
  $|\Delta\mathbf{h} \cdot \hat{\mathbf{v}}_{\text{probe}}| /
  \|\Delta\mathbf{h}\|$.
\end{enumerate}

\subsection{Evaluation metrics for decoder outputs}

For each generated continuation $g$, we compute: \textbf{Hedge score}
$H(g)$ (count of 28 standard clinical hedging markers;
Appendix~\ref{app:hedge_lexicon}), \textbf{Lexical diversity} $D(g)$
(alphabetic type-token ratio), \textbf{Token validity} (fraction of
tokens matching the NLTK English vocabulary), and \textbf{Perplexity}
$\mathrm{PPL}(g)$ under unsteered BioGPT to quantify distributional
shift (Appendix~\ref{app:eval_metrics}).

%% -----------------------------------------------------------------------
\section{Results}
%% -----------------------------------------------------------------------

\subsection{Pooling strategy determines observable structure}
\label{sec:pooling}

Before any multi-model analysis, we confirm for BioGPT that the choice
of representation extraction method materially affects what structure is
visible. At layer 16, using mean-difference vectors on held-out BioScope
sentences, \textbf{last-token pooling} yields Cohen's $d = 0.088$
(held-out classification accuracy $57\%$), whereas \textbf{mean pooling}
yields Cohen's $d = 1.331$ (accuracy $73\%$). The 15-fold improvement in
effect size reflects the fact that last-token hidden states in biomedical
sentences are dominated by the representation of the final content token
(often a medical noun such as \emph{pneumonia} or \emph{atelectasis}).
This representation is highly informative about the sentence's topic but
carries little signal about the speaker's epistemic stance. Mean pooling
averages across token positions, surfacing sentence-level properties such
as hedging that are distributed across the sequence. All subsequent
analysis uses mean-pooled hidden states.

\subsection{Uncertainty is linearly decodable across all seven models}
\label{sec:probing}

Table~\ref{tab:probe_all} reports probe accuracy at the anchor and
selected steering layers for all seven models. Three cross-model patterns
are notable.

\begin{table}[ht]
\centering
\caption{Linear probe CV accuracy for uncertainty classification. Layer
  $\hat{\ell}$ is the window-selected steering layer. BioGPT-Large and
  BioMedLM results are based on preliminary experimental runs.}
\label{tab:probe_all}
\begin{tabular}{lrrrrr}
\toprule
\textbf{Model} & $N$ & $\hat{\ell}$ & \textbf{L0 acc} &
  \textbf{Acc@}$\hat{\ell}$ & \textbf{Std@}$\hat{\ell}$ \\
\midrule
BioGPT          & 24 & 16 & 82.00\% & 85.50\% & 2.44\% \\
BioGPT-Large    & 48 & 31 & 82.50\% & 87.25\% & 1.63\% \\
BioMedLM        & 32 & 21 & 55.25\% & 86.00\% & --     \\
\midrule
BioBERT         & 12 &  8 & 78.00\% & 86.00\% & 3.24\% \\
ClinicalBERT    & 12 &  6 & 77.00\% & 84.75\% & 2.98\% \\
BlueBERT        & 12 & 10 & 77.25\% & 83.50\% & 5.18\% \\
SciBERT         & 12 & 10 & 80.00\% & 86.25\% & 4.76\% \\
\bottomrule
\end{tabular}
\end{table}

\paragraph{Finding 1: Uncertainty is robustly decodable in all models.}
All seven models achieve between 83\% and 87\% CV accuracy at the
selected steering layer. The range is narrow (4pp), indicating that
strong linear encoding of uncertainty is a general property of this
class of biomedical language models, not specific to any particular
architecture or pretraining distribution.

\paragraph{Finding 2: Token-embedding uncertainty signal depends on model family.}
Layer 0 accuracy---which measures how much uncertainty signal is already
present in the token embeddings before any transformer processing---
differs markedly across models. BioGPT achieves 82.00\% and BioGPT-Large
82.50\% at layer 0, indicating that a large fraction of the uncertainty
signal is available from surface lexical features alone. In contrast,
BioMedLM achieves only 55.25\% at layer 0 (near chance), suggesting that
for this model, uncertainty is not present in the token embeddings but is
built entirely through transformer computation. The encoder models
(BioBERT family) occupy an intermediate position (77--80\%), consistent
with bidirectional attention enabling richer embedding-level
contextualization even in early layers.

\paragraph{Finding 3: Probing curves are uniformly strong at middle-to-late layers.}
For all models, the accuracy curve rises through the first third of the
network and plateaus through the middle and upper two-thirds, with
minimal degradation in the final layers. This plateau structure is the
empirical motivation for the ±2 window in our layer selection procedure:
within a neighbourhood of any point in the upper two-thirds of the
network, accuracy differences are small (typically $<2$pp).

Full layer-by-layer probing tables for all models are provided in
Appendix~\ref{app:probe_results}.

\subsection{The uncertainty representation is effectively one-dimensional}

To test whether the $\sim$86\% accuracy reflects a rich multi-dimensional
uncertainty representation or a single dominant direction, we compare a
\textbf{full-dimensional probe} (logistic regression on the full hidden
state) against a \textbf{1D projection probe} (logistic regression on the
scalar projection onto $\mathbf{v}_{\mathrm{probe}}$), both under nested
5-fold cross-validation. For BioGPT at layer 16, the full-dimensional
probe achieves $85.50\%$ and the 1D projection achieves $85.50\%$---
identical within noise (standard deviations $2.44\%$ and $2.32\%$
respectively). A single linear direction captures essentially all of the
classifier-accessible uncertainty information. This finding holds across
encoder models as well: at each model's selected layer, the 1D projection
retains $>98\%$ of the full-dimensional probe accuracy.

\subsection{The representation survives surface-feature orthogonalization}

For BioGPT, we verify that the 85.50\% probing accuracy is not driven by
trivial surface correlates. We compute three projection classifiers:
(1)~the raw probe direction; (2)~the length-orthogonalized direction; and
(3)~the fully length-and-hedge-orthogonalized direction. Accuracies at
layer 16 are $85.50\%$, $85.75\%$, and $85.00\%$ respectively---
essentially identical. Removing both obvious surface confounds reduces
accuracy by less than one percentage point. We conclude that the linear
uncertainty encoding is not reducible to sentence length or
hedge-vocabulary frequency.

\subsection{Decoder steering: calibrated productive range}
\label{sec:decsteer}

Figure~\ref{fig:tradeoff} plots the four evaluation metrics for BioGPT
as functions of $\alpha_{\mathrm{frac}}$, using dynamically calibrated
hidden norms ($\|\mathbf{h}_{16}\|_{\text{est}} \approx 205$). Four
findings emerge.

\begin{figure}[ht]
\centering
\includegraphics[width=0.95\textwidth]{./plots/fig_steering_tradeoff_v2.png}
\caption{Four-metric steering characterization at BioGPT layer 16 as a
function of $\alpha_{\mathrm{frac}}$, comparing the probe direction
(solid circles) and orthogonalized direction (dashed squares). Dynamic
hidden-norm calibration is used throughout.}
\label{fig:tradeoff}
\end{figure}

\paragraph{Finding 1: A genuine steering effect exists at low magnitudes.}
At $\alpha_{\mathrm{frac}} = 0.05$ with the probe direction, hedge score
increases from the $0.38$ baseline to $0.98$ ($2.6\times$) while lexical
diversity ($0.95$) and perplexity ($20.2$) remain essentially at baseline
values. This is the signature of a clean semantic intervention: the
target property increases while other generation quality measures are
preserved.

\paragraph{Finding 2: Dynamic calibration extends the productive range.}
With statically estimated norms ($\|\mathbf{h}\|_{\text{est}} = 616$),
perplexity at $\alpha_{\mathrm{frac}} = 0.25$ exceeds 400---a collapse
by any measure. With dynamically calibrated norms ($\approx 205$), the
same $\alpha_{\mathrm{frac}}$ corresponds to a 3$\times$ smaller absolute
perturbation; perplexity remains below 45 and lexical diversity above
0.85. This concretely demonstrates that accurate norm calibration is not
a technical detail but determines whether a given $\alpha_{\mathrm{frac}}$
falls in the productive range or the collapse regime.

\paragraph{Finding 3: Hedge score is non-monotonic in $\alpha_{\mathrm{frac}}$.}
The probe direction's hedge score decreases between
$\alpha_{\mathrm{frac}} = 0.05$ ($H = 0.98$) and
$\alpha_{\mathrm{frac}} = 0.10$ ($H = 0.76$). This non-monotonicity is
diagnostic: the high hedge scores at large $\alpha_{\mathrm{frac}}$
represent a collapse into a degenerate register where hedging vocabulary
dominates at the expense of semantic content.

\paragraph{Finding 4: Orthogonalization does not substantially extend the productive range.}
Comparing the two directions at matched $\alpha_{\mathrm{frac}}$ reveals
small quantitative differences but no qualitative rescue. Both directions
show similar onset of perplexity increase, and both produce similar
generation quality at collapse. This is consistent with the interpretation
that collapse is a property of large activation perturbations in general,
not of the specific surface features encoded in the probe direction.

\subsection{Encoder representation steering}
\label{sec:encsteer}

Table~\ref{tab:enc_steer} reports the five encoder steering metrics at
$\alpha_{\mathrm{frac}} \in \{0, 0.05, 0.10, 0.15\}$ for the probe
direction, averaged across the 80-sentence test set. All models show
consistent, positive $\Delta_{\text{proj}}$ and $\Delta_{\text{prob}}$
at non-zero $\alpha_{\mathrm{frac}}$, indicating that the steering vector
predictably shifts representations toward the uncertain pole.

Several observations are worth noting. First, the on-axis fraction is
high ($>0.95$) at small $\alpha_{\mathrm{frac}}$, indicating that the
injected perturbation lands primarily on the probe direction with little
``leakage'' into orthogonal dimensions. Second, the flip rate increases
monotonically with $\alpha_{\mathrm{frac}}$, consistent with
representations crossing the probe's decision boundary as perturbation
magnitude grows. Third, across the four encoder models, the
representation shifts are consistent in direction but differ in
magnitude---ClinicalBERT and BlueBERT show smaller shifts than BioBERT
and SciBERT, consistent with their somewhat lower probe accuracy at the
selected layer.

\begin{table}[ht]
\centering
\caption{Encoder representation steering metrics (probe direction) at
  selected layers, averaged over 80 test sentences. Results from new
  5-metric pipeline.}
\label{tab:enc_steer}
\begin{tabular}{llrrrrr}
\toprule
\textbf{Model} & $\alpha_{\text{frac}}$ &
  $\Delta_{\text{proj}}$ & $\Delta_{\text{prob}}$ &
  \text{flip\%} & $\Delta_{\text{cos}}$ & \text{on-axis} \\
\midrule
\multirow{4}{*}{BioBERT (L8)}
  & 0.00 & +0.000 & +0.000 &  0.0\% & +0.000 & -- \\
  & 0.05 & +0.321 & +0.083 &  8.8\% & +0.012 & 0.971 \\
  & 0.10 & +0.643 & +0.167 & 22.5\% & +0.024 & 0.969 \\
  & 0.15 & +0.964 & +0.251 & 38.8\% & +0.036 & 0.967 \\
\midrule
\multirow{4}{*}{ClinicalBERT (L6)}
  & 0.00 & +0.000 & +0.000 &  0.0\% & +0.000 & -- \\
  & 0.05 & +0.275 & +0.069 &  6.3\% & +0.009 & 0.968 \\
  & 0.10 & +0.550 & +0.139 & 16.3\% & +0.019 & 0.966 \\
  & 0.15 & +0.825 & +0.209 & 28.8\% & +0.028 & 0.964 \\
\midrule
\multirow{4}{*}{BlueBERT (L10)}
  & 0.00 & +0.000 & +0.000 &  0.0\% & +0.000 & -- \\
  & 0.05 & +0.289 & +0.072 &  7.5\% & +0.010 & 0.969 \\
  & 0.10 & +0.578 & +0.145 & 18.8\% & +0.021 & 0.967 \\
  & 0.15 & +0.866 & +0.217 & 32.5\% & +0.031 & 0.965 \\
\midrule
\multirow{4}{*}{SciBERT (L10)}
  & 0.00 & +0.000 & +0.000 &  0.0\% & +0.000 & -- \\
  & 0.05 & +0.334 & +0.086 &  8.8\% & +0.013 & 0.973 \\
  & 0.10 & +0.668 & +0.172 & 23.8\% & +0.026 & 0.971 \\
  & 0.15 & +1.002 & +0.258 & 40.0\% & +0.039 & 0.969 \\
\bottomrule
\end{tabular}
\end{table}

\subsection{Summary of findings}

Across seven models and two architectural families, uncertainty is
consistently linearly encodable from mean-pooled hidden states at
middle-to-late layers, with CV accuracy ranging from 83--87\%. The
representation is effectively one-dimensional and survives
orthogonalization against length and hedge-vocabulary confounds. The
2/3-depth layer selection rule is empirically validated by probe data,
and the $\pm 2$ window refinement yields a meaningful improvement for at
least one model (ClinicalBERT: $+3.25$pp). For decoder models, activation
steering works at low magnitudes, but the productive range depends
critically on accurate hidden-norm calibration. For encoder models, five
complementary metrics show consistent, targeted representation shifts even
at small magnitudes.

%% -----------------------------------------------------------------------
\section{Discussion \& Conclusion}
%% -----------------------------------------------------------------------

Our results tell a consistent cross-model story.

\paragraph{Uncertainty is cleanly decodable across all seven models.}
The uniformity of probe accuracy (83--87\%) across models with different
architectures, sizes, and pretraining corpora is striking. It suggests
that the capacity to linearly encode epistemic stance is not a property
of any specific checkpoint or architecture, but a general feature of
biomedical language models trained on scientific or clinical text where
hedged language is prevalent. The single most informative cross-model
contrast is the layer-0 accuracy: BioGPT and BioGPT-Large start at
82\%, indicating that their uncertainty signal is largely available in
token embeddings (i.e., in the lexical statistics of the pretraining
vocabulary), whereas BioMedLM starts at 55\% (near chance), indicating
that its uncertainty encoding is built entirely by transformer
computation rather than embedded in the token representation.

\paragraph{Layer selection is not arbitrary.}
The 2/3-depth rule provides a principled, architecture-invariant
starting point for layer selection, and the $\pm 2$ window ensures that
locally better layers within the same semantic processing stage are not
missed. The ClinicalBERT case---where the L8 anchor achieves 81.50\%
but the window correctly identifies L6 at 84.75\% ($+3.25$pp) with lower
variance---illustrates that the window refinement catches real
differences. The window's restriction to $\pm 2$ is deliberate: it
prevents selection of layers far from the anchor that may be better
probes but represent a different processing stage (e.g., L5 for BioGPT
achieves 86.25\% globally but falls outside the upper-two-thirds window).

\paragraph{Linear decodability does not imply clean steerability.}
For decoder models, the productive steering range is narrow: at low
$\alpha_{\mathrm{frac}}$ the target property increases cleanly, but
beyond a model-specific threshold, generation quality degrades. With
dynamically calibrated hidden norms, this threshold is higher than
previously reported---the earlier apparent collapse at
$\alpha_{\mathrm{frac}} = 0.25$ was a calibration artifact (effective
perturbation 3$\times$ too large). But even with accurate calibration,
the productive range remains bounded. This is consistent with the
theoretical interpretation that steering along any direction extracted
from in-distribution activations pushes the hidden state out of
distribution once the perturbation is large enough, regardless of which
surface features the direction happens to be correlated with.

For encoder models, we measure representation shifts rather than
generation quality. The five-metric suite shows consistent, targeted
shifts in all four models. High on-axis fractions ($>0.96$) confirm
that the perturbation lands precisely on the probe direction with
minimal side effects, at least within the representation space. Whether
these representation shifts translate to meaningful changes in downstream
use of encoder models (e.g., as feature extractors or classification
heads) is a natural next question not addressed here.

\paragraph{Multi-metric evaluation is necessary.}
Our final observation is methodological. Hedge score alone at large
$\alpha_{\mathrm{frac}}$ would give a misleadingly optimistic picture of
steering success: the high hedge counts at collapse come from
individually valid hedging tokens strung together outside any coherent
linguistic structure. The same logic applies to encoder steering: a
large $\Delta_{\text{proj}}$ without the corresponding probability shift
and flip rate would be difficult to interpret. For steering interventions
in safety-sensitive domains, we argue that multi-metric evaluation---
combining the target behavior metric with at least one fluency metric and
one diversity metric for decoder models, or with probability and
boundary-crossing metrics for encoder models---should be a minimum
standard, not an optional robustness check.

%% -----------------------------------------------------------------------
\section{Limitations}
%% -----------------------------------------------------------------------

\paragraph{Incomplete coverage of decoder models.}
While probing results are available for all seven models, the full
decoder steering sweep (BioGPT-Large, BioMedLM) is ongoing at the time
of writing. The steering analysis in Section~\ref{sec:decsteer} is
complete for BioGPT; results for the larger decoder models will be
incorporated as they become available.

\paragraph{Encoder generation is not available.}
For the four encoder models, we measure representation-level effects
only. Whether these representation shifts translate to behavioral changes
in downstream tasks (e.g., masked language modeling, classification heads)
remains an open question.

\paragraph{Contrast set size and composition.}
Our steering direction is built from 200 uncertain and 200 certain
BioScope sentences. BioScope itself mixes several genres (abstracts, full
papers, clinical notes) and several speculation types (epistemic hedging,
conditionals, hypotheticals); a narrower annotation scheme focused
specifically on clinical epistemic uncertainty might yield a sharper
representation.

\paragraph{Surface-feature orthogonalization is partial.}
We orthogonalize against sentence length and hedge-word count. Other
features---specific bigrams, syntactic structure, discourse markers---
remain uncontrolled. The claim that uncertainty ``is not reducible to
surface features'' should be read as ``not reducible to these two
surface features.''

\paragraph{Automatic hedging measurement.}
Our hedge score is a count of 28 lexicon terms. It detects the presence
of hedging markers but does not assess whether hedging is clinically
appropriate. A radiologist reading the steered outputs might judge
hedging as correct or misplaced; we do not measure this.

\paragraph{Window size is a free parameter.}
The $\pm 2$ window in our layer selection procedure is fixed and not
cross-validated. A window of $\pm 1$ or $\pm 3$ would select the same
layer for most models in our study, but the choice deserves explicit
acknowledgment. The window represents approximately one transformer block
in each direction, which we argue is the smallest granularity of
representational change; a wider window would risk including layers at
qualitatively different processing stages.

\paragraph{Single prompt domain and hardware.}
Decoder experiments use radiology-flavored prompts only. Steering
behavior may differ for other clinical genres. Experiments run on Apple
MPS (laptop experiments) and a single CUDA GPU (larger models); exact
reproducibility across hardware is not guaranteed.

%% -----------------------------------------------------------------------
\appendix
\section{Experimental Details}
\label{app:details}
%% -----------------------------------------------------------------------

\paragraph{Hyperparameters.}
Generation: temperature $0.8$, top-$p$ $0.9$, max new tokens $60$, min
new tokens $20$, repetition penalty $1.2$. Logistic regression: $C=0.1$,
max\_iter $2000$, random\_state $42$. 5-fold stratified CV, seed 42.
Sample sizes: 200 uncertain + 200 certain BioScope sentences for probe
training; 80 balanced sentences (held out from training) for encoder
steering evaluation; 20 seeds $\times$ 5 prompts $=$ 100 generations per
decoder configuration.

\paragraph{Reproducibility.}
All experiments use fixed random seeds. BioScope XML files are
enumerated in alphabetical order before sampling, ensuring cross-platform
consistency of the 400-sentence contrast set. Per-sample generations,
hyperparameters, and environment metadata (platform, Python, PyTorch,
Transformers versions) are saved with each results file. Code will be
released upon acceptance.

\paragraph{Compute.}
Small models (BioGPT, all BERT variants) run on a single Apple M-series
laptop (MPS, 16GB unified memory). BioGPT-Large and BioMedLM run on a
single CUDA GPU. All models use float32 precision; quantization was not
applied to any model in this study.

%% -----------------------------------------------------------------------
\section{Hedge Marker Lexicon}
\label{app:hedge_lexicon}
%% -----------------------------------------------------------------------

To compute the hedge score $H(g)$, we count occurrences of the following
28 lexical markers, drawn from the BioScope speculation cue vocabulary
and established clinical hedging conventions:
\emph{may, might, could, would, should; suggest, suggests, suggestive;
possible, possibly, possibility; probable, probably, probability;
likely, unlikely; appear, appears, appearing; seem, seems; suspect,
suspected, suspicion; consistent with, cannot exclude, rule out}.

%% -----------------------------------------------------------------------
\section{Full Probing Results}
\label{app:probe_results}
%% -----------------------------------------------------------------------

Tables~\ref{tab:probe_biogpt}--\ref{tab:probe_encoders} provide
layer-by-layer 5-fold CV probe accuracies for all models with complete
probing results. The anchor layer ($\ell^* = \mathrm{round}(2N/3)$) and
window-selected layer ($\hat{\ell}$) are indicated.

\begin{table}[ht]
\centering
\caption{BioGPT linear probe accuracy, all 25 layers. Anchor $\ell^*=16$,
  selected $\hat{\ell}=16$.}
\label{tab:probe_biogpt}
\small
\begin{tabular}{rrrr}
\toprule
Layer & CV Acc & Std & 95\% CI \\
\midrule
0  & 82.00\% & 4.38\% & [78.16\%, 85.84\%] \\
1  & 84.25\% & 4.73\% & [80.11\%, 88.39\%] \\
2  & 84.50\% & 3.26\% & [81.64\%, 87.36\%] \\
3  & 84.25\% & 3.01\% & [81.61\%, 86.89\%] \\
4  & 84.75\% & 2.24\% & [82.79\%, 86.71\%] \\
5  & 86.25\% & 1.98\% & [84.52\%, 87.98\%] \\
6  & 85.50\% & 2.09\% & [83.67\%, 87.33\%] \\
7  & 84.25\% & 3.38\% & [81.29\%, 87.21\%] \\
8  & 84.00\% & 2.40\% & [81.89\%, 86.11\%] \\
9  & 85.50\% & 2.59\% & [83.23\%, 87.77\%] \\
10 & 84.75\% & 2.40\% & [82.64\%, 86.86\%] \\
11 & 84.50\% & 2.44\% & [82.36\%, 86.64\%] \\
12 & 86.00\% & 1.85\% & [84.37\%, 87.63\%] \\
13 & 84.75\% & 1.63\% & [83.32\%, 86.18\%] \\
14 & 84.75\% & 2.05\% & [82.95\%, 86.55\%] \\
15 & 85.00\% & 1.98\% & [83.27\%, 86.73\%] \\
\textbf{16$^*\hat{}$} & \textbf{85.50\%} & \textbf{2.44\%} & \textbf{[83.36\%, 87.64\%]} \\
17 & 84.75\% & 2.24\% & [82.79\%, 86.71\%] \\
18 & 85.25\% & 2.85\% & [82.75\%, 87.75\%] \\
19 & 85.50\% & 2.44\% & [83.36\%, 87.64\%] \\
20 & 84.75\% & 3.11\% & [82.02\%, 87.48\%] \\
21 & 84.25\% & 2.88\% & [81.73\%, 86.77\%] \\
22 & 82.25\% & 2.24\% & [80.29\%, 84.21\%] \\
23 & 81.50\% & 2.24\% & [79.54\%, 83.46\%] \\
24 & 83.25\% & 1.90\% & [81.59\%, 84.91\%] \\
\bottomrule
\end{tabular}
\end{table}

\begin{table}[ht]
\centering
\caption{BERT encoder models: linear probe accuracy at all 13 layers
  (L0 = embedding, L1--L12 = transformer blocks).
  Anchor $\ell^*=8$ for all; window-selected $\hat{\ell}$ indicated.}
\label{tab:probe_encoders}
\small
\begin{tabular}{rrrrrrrrrrrrrr}
\toprule
 & \multicolumn{3}{c}{\textbf{BioBERT} ($\hat{\ell}=8$)}
 & \multicolumn{3}{c}{\textbf{ClinicalBERT} ($\hat{\ell}=6$)}
 & \multicolumn{3}{c}{\textbf{BlueBERT} ($\hat{\ell}=10$)}
 & \multicolumn{3}{c}{\textbf{SciBERT} ($\hat{\ell}=10$)} \\
\cmidrule(lr){2-4}\cmidrule(lr){5-7}\cmidrule(lr){8-10}\cmidrule(lr){11-13}
L & Acc & Std & & Acc & Std & & Acc & Std & & Acc & Std & \\
\midrule
0  & 78.0\% & 2.3\% && 77.0\% & 4.6\% && 77.3\% & 4.9\% && 80.0\% & 3.9\% \\
1  & 80.8\% & 2.9\% && 77.5\% & 2.3\% && 78.5\% & 6.3\% && 82.5\% & 4.7\% \\
2  & 81.0\% & 3.1\% && 81.8\% & 4.4\% && 78.0\% & 5.7\% && 84.0\% & 3.5\% \\
3  & 85.3\% & 3.5\% && 83.8\% & 4.4\% && 80.0\% & 3.6\% && 83.5\% & 3.2\% \\
4  & 83.5\% & 3.9\% && 82.8\% & 6.6\% && 81.8\% & 3.7\% && 84.8\% & 3.7\% \\
5  & 83.8\% & 3.8\% && 83.0\% & 4.5\% && 78.8\% & 6.1\% && 84.8\% & 4.1\% \\
6  & 84.3\% & 3.3\% && \textbf{84.8\%}$^\wedge$ & \textbf{3.0\%} && 80.8\% & 4.1\% && 84.0\% & 5.3\% \\
7  & 84.3\% & 3.4\% && 81.0\% & 4.1\% && 80.8\% & 3.4\% && 85.8\% & 4.8\% \\
8$^*$& \textbf{86.0\%}$^\wedge$ & \textbf{3.2\%} && 81.5\% & 4.5\% && 83.3\% & 3.9\% && \textbf{86.0\%} & 3.7\% \\
9  & 85.8\% & 3.6\% && 81.8\% & 3.9\% && 83.5\% & 5.6\% && 85.0\% & 5.0\% \\
10 & 84.5\% & 4.6\% && 82.0\% & 2.1\% && \textbf{83.5\%}$^\wedge$ & \textbf{5.2\%} && \textbf{86.3\%}$^\wedge$ & \textbf{4.8\%} \\
11 & 84.0\% & 5.2\% && 82.8\% & 1.9\% && 82.5\% & 6.2\% && 86.0\% & 3.9\% \\
12 & 80.8\% & 4.7\% && 81.3\% & 2.8\% && 81.5\% & 4.1\% && 84.5\% & 5.1\% \\
\bottomrule
\multicolumn{13}{l}{\small $^*$ = 2/3 anchor; $^\wedge$ = window-selected layer.}
\end{tabular}
\end{table}

%% -----------------------------------------------------------------------
\section{Formal Evaluation Metrics}
\label{app:eval_metrics}
%% -----------------------------------------------------------------------

\paragraph{Lexical diversity.}
$D(g)$: type-token ratio over alphabetic words:
\begin{equation}
D(g) = \frac{|\text{unique words}(g)|}{|\text{words}(g)|}
\end{equation}

\paragraph{Perplexity.}
$\mathrm{PPL}(g)$: computed under vanilla (unsteered) BioGPT to measure
distributional shift:
\begin{equation}
\mathrm{PPL}(g) = \exp\!\left(-\frac{1}{|g|}
\sum_{t=1}^{|g|}\log p_\theta(g_t \mid g_{<t})\right)
\end{equation}
where $p_\theta$ is BioGPT without any steering hooks applied.

%% -----------------------------------------------------------------------
\section{Detailed Steering Results (BioGPT)}
\label{app:steering_table}
%% -----------------------------------------------------------------------

Table~\ref{tab:steering_app} provides numerical values for the BioGPT
decoder steering sweep using dynamically calibrated hidden norms
($\|\mathbf{h}_{16}\|_{\text{est}} \approx 205$).

\begin{table}[ht]
\centering
\caption{BioGPT decoder steering metrics at layer 16, calibrated norms,
  averaged over 100 samples per configuration. Bold values mark peak
  hedge score while retaining near-baseline lexical diversity.}
\label{tab:steering_app}
\begin{tabular}{lrrrr}
\toprule
$\alpha_{\mathrm{frac}}$ & Hedge & Lex.~div. & Tok.~val. & PPL \\
\midrule
\multicolumn{5}{l}{\textit{Probe direction}} \\
0.000 & 0.38 & 0.94 & 0.87 & 16.1 \\
0.025 & 0.74 & 0.94 & 0.87 & 18.0 \\
0.050 & \textbf{0.98} & \textbf{0.95} & 0.89 & 20.2 \\
0.075 & 0.92 & 0.94 & 0.92 & 28.4 \\
0.100 & 0.76 & 0.91 & 0.93 & 34.7 \\
0.125 & 0.89 & 0.88 & 0.94 & 38.1 \\
0.150 & 1.09 & 0.85 & 0.95 & 40.2 \\
0.200 & 1.82 & 0.79 & 0.95 & 42.3 \\
0.250 & 2.47 & 0.74 & 0.94 & 42.8 \\
\midrule
\multicolumn{5}{l}{\textit{Orthogonalized direction}} \\
0.000 & 0.38 & 0.94 & 0.87 & 16.1 \\
0.025 & 0.64 & 0.94 & 0.88 & 17.3 \\
0.050 & 0.74 & 0.94 & 0.89 & 21.4 \\
0.075 & \textbf{1.00} & \textbf{0.94} & 0.92 & 26.8 \\
0.100 & 0.62 & 0.92 & 0.94 & 31.5 \\
0.125 & 0.71 & 0.90 & 0.95 & 35.2 \\
0.150 & 0.54 & 0.88 & 0.96 & 37.9 \\
0.200 & 1.19 & 0.81 & 0.96 & 40.1 \\
0.250 & 1.97 & 0.76 & 0.95 & 41.6 \\
\bottomrule
\end{tabular}
\end{table}

%% -----------------------------------------------------------------------
%% Bibliography (add your .bib file)
%% -----------------------------------------------------------------------
\bibliographystyle{plainnat}
\bibliography{references}

\end{document}
